%%=============================================================================
%% uitleg blazor
%%=============================================================================


\chapter*{Wat is blazor?}

Blazor is een framework van Microsoft waarmee je webapplicaties kan bouwen 
met C\# in plaats van JavaScript. Dit maakt het interessant voor bedrijven 
die al werken met het .NET-ecosysteem, zoals Victor Buyck Steel Construction.

\bigskip
Blazor bestaat in twee varianten. De eerste is \textbf{Blazor Server}, waarbij 
de applicatie volledig op de server draait en de browser enkel de UI toont. 
De communicatie tussen browser en server verloopt via SignalR. Dit betekent 
dat er altijd een actieve verbinding nodig is met de server.

\bigskip
De tweede variant is \textbf{Blazor WebAssembly}. Hierbij wordt de volledige 
applicatie gedownload naar de browser en lokaal uitgevoerd via WebAssembly. 
WebAssembly is een standaard die moderne browsers ondersteunen en die het 
mogelijk maakt om code te draaien aan near-native snelheid, zonder 
serververbinding. Een bijkomend voordeel is dat de applicatie als een 
standalone app kan draaien, wat ideaal is voor een onthaaltoestel dat 
niet altijd verbonden hoeft te zijn.

